\documentclass[9pt,landscape,a4paper]{extarticle}
\usepackage[utf8]{inputenc}
\usepackage[vietnamese]{babel}
\usepackage[T1]{fontenc}
\usepackage{fontspec}
\setmainfont{Libertinus Serif} % As per user's font

% Layout and multicolumn
\usepackage[margin=0.8cm, top=1.5cm, bottom=1.5cm]{geometry}
\usepackage{multicol}
\setlength{\columnsep}{0.8cm}
\setlength{\columnseprule}{0.1pt}

% Math and symbols
\usepackage{amsmath,amssymb}
\usepackage{wasysym}
\usepackage{xcolor}

% Custom Colors (Stanford Cheatsheet style)
\definecolor{stanfordblue}{RGB}{1, 104, 149}
\definecolor{stanfordred}{RGB}{140, 21, 21}
\definecolor{stanfordgreen}{RGB}{0, 100, 0}
\definecolor{cyanbox}{RGB}{220, 245, 255}

% Section styling (mimicking the sheet)
\usepackage{titlesec}
\titleformat{\section}
  {\normalfont\normalsize\bfseries}{\thesection}{1em}{}
\titlespacing*{\section}{0pt}{1.5ex plus 1ex minus .2ex}{0.5ex plus .2ex}

\titleformat{\subsection}
  {\normalfont\small\bfseries}{\thesubsection}{1em}{}
\titlespacing*{\subsection}{0pt}{1ex plus 1ex minus .2ex}{0.5ex plus .2ex}

% Fancy header
\usepackage{fancyhdr}
\pagestyle{fancy}
\fancyhf{}
\renewcommand{\headrulewidth}{0.4pt}
\renewcommand{\footrulewidth}{0.4pt}

% Header content
\lhead{\textsc{MAT3562E} -- \textsc{COMPUTER VISION}}
\rhead{\colorbox{cyan!20}{\makebox[4cm][c]{\textbf{Bùi Quang Chiến -- 23001837}}}}

% Footer content
\lfoot{\textsc{HUS -- VNU}}
\cfoot{\thepage}
\rfoot{\textsc{SPRING 2026}}

% Dense lists
\usepackage{enumitem}
\setlist{nosep, leftmargin=1.2em}
\setlist[itemize,1]{label={\textcolor{stanfordgreen}{\footnotesize$\square$}}}
\setlist[itemize,2]{label=--}

\begin{document}
\begin{center}
    {\Large \bfseries VIP Cheatsheet: LAB04_05} \\
    \vspace{0.3cm}
    {\large \textsc{Bùi Quang Chiến}} \\
    \vspace{0.2cm}
    {\small April 2026}
\end{center}
\vspace{0.3cm}

\begin{multicols*}{3}

\begin{center}
    {\Large \bfseries VIP Cheatsheet: LAB04_05} \\
    \vspace{0.3cm}
    {\large \textsc{Bùi Quang Chiến}} \\
    \vspace{0.2cm}
    {\small April 2026}
\end{center}
\vspace{0.3cm}

\begin{multicols*}{3}

\begin{center}
  {\large\bfseries\color{hdr} LAB 04-05 CHEATSHEET --- Biến đổi Hình \& Hình thái học} \\[1pt]
  {\small Bùi Quang Chiến \textbullet{} 23001837 \textbullet{} MAT3562E Computer Vision 2025--2026}
\end{center}
\hrule\vspace{2pt}

\begin{multicols}{3}

\section{1. Biến đổi Hình học (Geometric Transformations)}

\subsection*{Tịnh tiến (Translation)}
Ma trận tịnh tiến: $M = \begin{bmatrix} 1 & 0 & t_x \\ 0 & 1 & t_y \end{bmatrix}$.\\
\begin{lstlisting}[language=Python]
cv2.warpAffine(img, M, (w, h))
\end{lstlisting}

\subsection*{Co giãn (Scaling)}
Ma trận co giãn: $M = \begin{bmatrix} s_x & 0 & 0 \\ 0 & s_y & 0 \end{bmatrix}$.\\
\begin{lstlisting}[language=Python]
cv2.resize(img, None, fx=s_x, fy=s_y, interpolation=cv2.INTER_LINEAR)
\end{lstlisting}

\subsection*{Xoay (Rotation)}
Xoay quanh tâm \texttt{center} với góc \texttt{angle}:
\begin{lstlisting}[language=Python]
M = cv2.getRotationMatrix2D(center, angle, scale)
cv2.warpAffine(img, M, (new_w, new_h))
\end{lstlisting}

\subsection*{Kéo lệch (Shear / Skew)}
Ngang: $M = \begin{bmatrix} 1 & k & 0 \\ 0 & 1 & 0 \end{bmatrix}$. Dọc: $M = \begin{bmatrix} 1 & 0 & 0 \\ k & 1 & 0 \end{bmatrix}$.
\begin{lstlisting}[language=Python]
cv2.warpAffine(img, M, (new_w, new_h))
\end{lstlisting}

\begin{tipbox}
\textbf{Kiểm soát Canvas \& Padding:}\\
Dùng \texttt{borderMode=cv2.BORDER\_CONSTANT} và \texttt{borderValue=0} để điền viền đen. Ánh xạ ngược (Inverse Mapping) loại bỏ hiện tượng lỗ hở (holes).
\end{tipbox}

\section{2. Hình thái học (Morphological Operations)}

\begin{lstlisting}[language=Python]
# Hạt nhân cấu trúc (kernel)
kernel = cv2.getStructuringElement(cv2.MORPH_ELLIPSE, (5, 5))
\end{lstlisting}

\begin{infobox}
\textbf{Co (Erosion):} \texttt{cv2.erode(img, kernel)}\\
Làm mỏng đối tượng, xóa nét gai nhỏ, xẻ các đoạn dính hẹp trên đường biên. (Min cục bộ).
\end{infobox}

\begin{infobox}
\textbf{Giãn (Dilation):} \texttt{cv2.dilate(img, kernel)}\\
Làm dày đường nét đối tượng, nối đứt gãy bên trong đối tượng. (Max cục bộ).
\end{infobox}

\begin{infobox}
\textbf{Mở (Opening):} \texttt{cv2.morphologyEx(..., cv2.MORPH\_OPEN)}\\
\textbf{Co $\to$ Giãn.} Gạt bỏ sạn sáng nhỏ li ti bên ngoài vật thể, dọn sạch gai biên nhọn. Tỷ lệ cốt lõi được bảo toàn.
\end{infobox}

\begin{infobox}
\textbf{Đóng (Closing):} \texttt{cv2.morphologyEx(..., cv2.MORPH\_CLOSE)}\\
\textbf{Giãn $\to$ Co.} Trám đầy kẽ nứt hở hay lõm thụt sẹo li ti nằm trọn phía trong cấu trúc mạch.
\end{infobox}

\begin{notebox}
\textbf{Pipeline Y Tế Phổ Biến:}\\
Khử mờ $\to$ Threshold (Otsu) $\to$ \textbf{Opening} (xóa nhiễu rác ngoài biên) $\to$ \textbf{Closing} (vá thủng bên trong khối).
\end{notebox}

\end{multicols}

\end{multicols*}

\end{multicols*}
\end{document}
